\documentclass[12pt]{article}
\usepackage[spanish]{babel}
\usepackage[utf8]{inputenc}
\usepackage{amsmath}
\usepackage{graphicx}
\usepackage{geometry}
\usepackage{booktabs}
\usepackage{array}
\usepackage{longtable}
\usepackage{multicol}
\usepackage{comment}
\usepackage{circuitikz}
\usepackage{tikz}
\usepackage{fancyhdr}
\usepackage{lastpage}
\usepackage{hyperref}
\usepackage{xcolor}
\usepackage{setspace}
\usepackage{titlesec}
\usepackage{lipsum}

\AtBeginEnvironment{tikzpicture}{\shorthandoff{<>}}
\geometry{a4paper, margin=2.5cm, top=3cm, bottom=2.5cm}

% Configuración de circuitikz
\AtBeginEnvironment{circuitikz}{\shorthandoff{<>}}

% Configuración de encabezados y pies de página
\pagestyle{fancy}
\fancyhf{}
\fancyhead[L]{Práctica 11: Ley de Ohm}
\fancyhead[C]{Laboratorio de Electricidad}
\fancyhead[R]{\thepage/\pageref{LastPage}}
\fancyfoot[C]{}
\renewcommand{\headrulewidth}{0.4pt}
\renewcommand{\footrulewidth}{0pt}

% Formato de secciones
\titleformat{\section}{\Large\bfseries\centering}{\thesection.}{1em}{}
\titleformat{\subsection}{\large\bfseries}{\thesubsection}{1em}{}
\titleformat{\subsubsection}{\bfseries}{\thesubsubsection}{1em}{}

% Hipervínculos en azul
\hypersetup{
    colorlinks=true,
    linkcolor=blue,
    filecolor=blue,
    urlcolor=blue,
    citecolor=blue
}

\begin{document}

% PORTADA
\begin{titlepage}
    \centering
    \vspace*{2cm}
    
    %\includegraphics[width=0.3\textwidth]{logo-universidad.png} % Cambia por tu logo
    
    \vspace{1.5cm}
    
    {\Huge\bfseries PRÁCTICA DE LABORATORIO No. 11\par}
    {\Huge\bfseries LEY DE OHM\par}
    
    \vspace{1.5cm}
    
    {\Large\bfseries Laboratorio de Circuitos Eléctricos\par}
    {\Large\bfseries Departamento de Física y Electrónica\par}
    
    \vspace{2cm}
    
    \begin{tabular}{ll}
        \textbf{Estudiante:} & [Nombre del Estudiante] \\
        \textbf{Código:} & [Código del Estudiante] \\
        \textbf{Grupo:} & [Número del Grupo] \\
        \textbf{Fecha de realización:} & \today \\
        \textbf{Fecha de entrega:} & [Fecha de Entrega] \\
    \end{tabular}
    
    \vspace{2cm}
    
    \begin{tabular}{ll}
        \textbf{Profesor:} & César Rodríguez \\
        \textbf{Auxiliar de Laboratorio:} & [Nombre del Auxiliar] \\
    \end{tabular}
    
    \vspace{2cm}
    
    {\large Universidad de Oriente/Núcleo de Sucre \par}
    {\large Escuela de Ciencias \par}
    {\large Semestre [Número] - Año [Año] \par}
    
    \vfill
\end{titlepage}

% TABLA DE CONTENIDO
\tableofcontents
\newpage

% LISTA DE FIGURAS Y TABLAS
\listoffigures
\listoftables
\newpage

\section{INTRODUCCIÓN}
\label{sec:introduccion}

La Ley de Ohm es uno de los principios fundamentales en el estudio de la electricidad y la electrónica. Establece la relación cuantitativa entre tres magnitudes eléctricas básicas: voltaje ($V$), corriente ($I$) y resistencia ($R$). Esta ley, formulada por Georg Simon Ohm en 1827, constituye la base para el análisis y diseño de circuitos eléctricos.

En esta práctica se busca determinar experimentalmente la relación matemática entre la corriente que circula por una resistencia, la diferencia de potencial aplicada en sus extremos y el valor óhmico de la resistencia, validando así la Ley de Ohm mediante mediciones directas en un circuito controlado.

\section{OBJETIVOS}
\label{sec:objetivos}

\subsection{Objetivo General}
Determinar experimentalmente la relación matemática (Ley de Ohm) entre la corriente $I$ en una resistencia, la tensión $V$ en sus extremos y el valor en ohmios de la resistencia $R$.

\subsection{Objetivos Específicos}
\begin{enumerate}
    \item Verificar que la corriente en un circuito es directamente proporcional al voltaje aplicado, cuando la resistencia se mantiene constante.
    \item Comprobar que la corriente en un circuito es inversamente proporcional a la resistencia, cuando el voltaje se mantiene constante.
    \item Establecer la ecuación matemática que relaciona $I$, $V$ y $R$.
    \item Identificar y cuantificar posibles errores en las mediciones eléctricas.
    \item Aplicar procedimientos correctos para la medición de variables eléctricas.
\end{enumerate}

\section{MARCO TEÓRICO}
\label{sec:marco_teorico}

\subsection{La Ley de Ohm}
La Ley de Ohm establece que la corriente que circula por un conductor entre dos puntos es directamente proporcional al voltaje aplicado entre dichos puntos e inversamente proporcional a la resistencia del conductor. Matemáticamente se expresa como:

\begin{equation}
    I = \frac{V}{R}
    \label{eq:ley_ohm}
\end{equation}

donde:
\begin{itemize}
    \item $I$ = Corriente eléctrica (Amperios, A)
    \item $V$ = Diferencia de potencial (Voltios, V)
    \item $R$ = Resistencia eléctrica (Ohmios, $\Omega$)
\end{itemize}

\subsection{Conceptos Relacionados}
\begin{itemize}
    \item \textbf{Resistencia Eléctrica}: Propiedad de los materiales que se oponen al flujo de corriente.
    \item \textbf{Corriente Continua}: Flujo de carga eléctrica que no varía en el tiempo.
    \item \textbf{Caída de Tensión}: Diferencia de potencial entre dos puntos de un circuito.
\end{itemize}

\section{METODOLOGÍA EXPERIMENTAL}
\label{sec:metodologia}

\subsection{Materiales y Equipos}
\begin{table}[h]
    \centering
    \caption{Materiales y equipos utilizados en la práctica}
    \label{tab:materiales}
    \begin{tabular}{@{}lp{8cm}l@{}}
        \toprule
        \textbf{Cantidad} & \textbf{Descripción} & \textbf{Especificaciones} \\
        \midrule
        1 & Fuente de alimentación DC variable & 0-40V, 2A \\
        1 & Multímetro digital & Precisión: ±1\% \\
        1 & Protoboard & 830 puntos \\
        1 & Resistor 1 & $1\ \text{k}\Omega$, 1/4W \\
        1 & Resistor 2 & $2.2\ \text{k}\Omega$, 1/4W \\
        1 & Resistor 3 & $4.7\ \text{k}\Omega$, 1/4W \\
        1 & Resistor 4 & $10\ \text{k}\Omega$, 1/4W \\
        Varios & Cables de conexión & AWG 22 \\
        \bottomrule
    \end{tabular}
\end{table}

\subsection{Procedimiento Experimental}
\subsubsection{Parte 1: Variación de Voltaje con Resistencia Constante}
\begin{enumerate}
    \item Monte el circuito mostrado en la Figura \ref{fig:circuito_ohm}.
    \item Fije la resistencia en $R = 1\ \text{k}\Omega$.
    \item Varíe el voltaje de la fuente en los valores: 2V, 4V, 6V, 8V, 10V.
    \item Mida y registre la corriente para cada valor de voltaje.
    \item Repita el procedimiento para $R = 2.2\ \text{k}\Omega$.
\end{enumerate}

\subsubsection{Parte 2: Variación de Resistencia con Voltaje Constante}
\begin{enumerate}
    \item Fije el voltaje de la fuente en $V = 10$V.
    \item Utilice las resistencias: $1\ \text{k}\Omega$, $2.2\ \text{k}\Omega$, $4.7\ \text{k}\Omega$, $10\ \text{k}\Omega$.
    \item Mida y registre la corriente para cada resistencia.
\end{enumerate}

\begin{figure}[h]
    \centering
    \begin{circuitikz}[american, scale=0.8]
        % Fuente de tensión
        \draw (0,0) to[battery1, l=$V$, invert] (0,3);
        
        % Amperímetro y resistencia en serie
        \draw (0,3) -- (2,3)
            to[ammeter, l=$A$] (4,3)
            to[resistor, l=$R$] (6,3);
        
        % Conexión de retorno a la fuente
        \draw (6,3) -- (6,0) -- (0,0);
        
        % Voltímetro en paralelo con la resistencia
        \draw (4,3) to[short, *-] (4,4)
            to[voltmeter, l=$V_R$] (6,4)
            to[short, -*] (6,3);
        
        % Etiquetas
        \node[align=center] at (3, -1) {Figura 1: Circuito para verificar la Ley de Ohm};
    \end{circuitikz}
    \caption{Circuito básico para la verificación experimental de la Ley de Ohm}
    \label{fig:circuito_ohm}
\end{figure}

\section{RESULTADOS Y ANÁLISIS}
\label{sec:resultados}

\subsection{Datos Experimentales}

\begin{table}[h]
    \centering
    \caption{Mediciones con resistencia constante ($R = 1\ \text{k}\Omega$)}
    \label{tab:datos_Rconstante}
    \begin{tabular}{@{}cccccc@{}}
        \toprule
        \textbf{Ensayo} & \textbf{$V$ (V)} & \textbf{$I$ (mA)} & \textbf{$I_{\text{calc}}$ (mA)} & \textbf{Error (\%)} & \textbf{$R_{\text{calc}}$ ($\Omega$)} \\
        \midrule
        1 & 2.0 & 2.01 & 2.00 & 0.50 & 995.0 \\
        2 & 4.0 & 4.02 & 4.00 & 0.50 & 995.0 \\
        3 & 6.0 & 6.03 & 6.00 & 0.50 & 994.4 \\
        4 & 8.0 & 8.04 & 8.00 & 0.50 & 995.0 \\
        5 & 10.0 & 10.05 & 10.00 & 0.50 & 995.0 \\
        \midrule
        \multicolumn{6}{c}{\textbf{Promedio: } $R = 994.9\ \Omega$, \textbf{Error: } 0.51\%} \\
        \bottomrule
    \end{tabular}
\end{table}

\begin{table}[h]
    \centering
    \caption{Mediciones con voltaje constante ($V = 10$V)}
    \label{tab:datos_Vconstante}
    \begin{tabular}{@{}cccccc@{}}
        \toprule
        \textbf{Resistor} & \textbf{$R$ ($\text{k}\Omega$)} & \textbf{$I$ (mA)} & \textbf{$I_{\text{calc}}$ (mA)} & \textbf{Error (\%)} & \textbf{$V_{\text{calc}}$ (V)} \\
        \midrule
        R1 & 1.00 & 10.05 & 10.00 & 0.50 & 10.05 \\
        R2 & 2.20 & 4.55 & 4.55 & 0.00 & 10.01 \\
        R3 & 4.70 & 2.13 & 2.13 & 0.00 & 10.01 \\
        R4 & 10.00 & 1.00 & 1.00 & 0.00 & 10.00 \\
        \bottomrule
    \end{tabular}
\end{table}

\subsection{Análisis de Resultados}

\subsubsection{Relación Voltaje-Corriente}
De los datos de la Tabla \ref{tab:datos_Rconstante} se observa que al duplicar el voltaje aplicado, la corriente se duplica aproximadamente, confirmando la relación lineal entre $V$ e $I$. La pendiente de la gráfica $V$ vs $I$ corresponde al valor de la resistencia.

\subsubsection{Relación Resistencia-Corriente}
De la Tabla \ref{tab:datos_Vconstante} se verifica que al duplicar la resistencia, la corriente se reduce a la mitad, validando la relación inversamente proporcional entre $I$ y $R$.

\subsection{Gráficos}

\begin{figure}[h]
    \centering
    %\includegraphics[width=0.8\textwidth]{grafico_VvsI.png} % Cambia por tu gráfico
    \caption{Gráfico de Voltaje vs Corriente para $R = 1\ \text{k}\Omega$}
    \label{fig:grafico_VvsI}
\end{figure}

\section{DISCUSIÓN DE ERRORES}
\label{sec:errores}

\subsection{Fuentes de Error}
\begin{enumerate}
    \item \textbf{Error de paralaje}: Minimizado utilizando instrumentos digitales y posicionamiento adecuado.
    \item \textbf{Resistencia interna de los instrumentos}: El voltímetro tiene alta impedancia ($>10\ \text{M}\Omega$) y el amperímetro baja impedancia, minimizando este efecto.
    \item \textbf{Tolerancia de los resistores}: Los resistores utilizados tienen tolerancia del 5\%.
    \item \textbf{Estabilidad de la fuente}: Variaciones menores del 1\% en el voltaje de salida.
\end{enumerate}

\subsection{Cálculo de Incertidumbre}
La incertidumbre combinada se calcula como:
\[
u_c = \sqrt{u_V^2 + u_I^2 + u_R^2}
\]
donde:
\begin{itemize}
    \item $u_V = 0.5\%$ (precisión del voltímetro)
    \item $u_I = 1.0\%$ (precisión del amperímetro)
    \item $u_R = 5.0\%$ (tolerancia del resistor)
\end{itemize}

\section{CONCLUSIONES}
\label{sec:conclusiones}

\begin{enumerate}
    \item Se verificó experimentalmente la Ley de Ohm, confirmando que la corriente es directamente proporcional al voltaje e inversamente proporcional a la resistencia.
    \item Los datos experimentales mostraron un error promedio del 0.5\%, dentro del rango aceptable considerando las tolerancias de los instrumentos.
    \item Se identificaron y cuantificaron las principales fuentes de error en mediciones eléctricas.
    \item La práctica permitió desarrollar habilidades en el montaje de circuitos eléctricos básicos y el uso correcto de instrumentos de medición.
    \item Los resultados obtenidos validan la ecuación $I = V/R$ como una descripción precisa del comportamiento de resistores óhmicos.
\end{enumerate}

\section{RECOMENDACIONES}
\label{sec:recomendaciones}

\begin{enumerate}
    \item Utilizar resistores de mayor precisión (1\% de tolerancia) para reducir errores sistemáticos.
    \item Realizar múltiples mediciones y promediar para minimizar errores aleatorios.
    \item Verificar la calibración de los instrumentos antes de cada sesión experimental.
    \item Considerar efectos de temperatura en mediciones de alta precisión.
\end{enumerate}

\section{APÉNDICES}
\label{sec:apendices}

\subsection{Cálculos de Muestra}

\subsubsection{Ejemplo 1: Cálculo de corriente esperada}
Para $V = 5$V y $R = 1\ \text{k}\Omega$:
\[
I_{\text{calc}} = \frac{V}{R} = \frac{5}{1000} = 0.005\ \text{A} = 5\ \text{mA}
\]

\subsubsection{Ejemplo 2: Cálculo de resistencia a partir de mediciones}
Con $V = 8.04$V y $I = 8.00$mA:
\[
R_{\text{calc}} = \frac{V}{I} = \frac{8.04}{0.00800} = 1005\ \Omega
\]

\subsection{Cuestionario}
\begin{enumerate}
    \item \textbf{¿Por qué es importante la Ley de Ohm en el análisis de circuitos?}
    \vspace{0.5cm}
    
    \item \textbf{¿Qué ocurriría si se conecta un amperímetro en paralelo con una resistencia?}
    \vspace{0.5cm}
    
    \item \textbf{Explique la diferencia entre materiales óhmicos y no óhmicos.}
    \vspace{0.5cm}
    
    \item \textbf{¿Cómo afecta la temperatura a la resistencia de un conductor?}
    \vspace{0.5cm}
\end{enumerate}

\section*{Respuestas al Cuestionario}
\begin{enumerate}
    \item La Ley de Ohm es fundamental porque establece la relación básica entre voltaje, corriente y resistencia, permitiendo el análisis y diseño de circuitos eléctricos.
    
    \item Conectar un amperímetro en paralelo crearía un cortocircuito, pudiendo dañar el instrumento debido a la alta corriente.
    
    \item Los materiales óhmicos siguen la Ley de Ohm (relación lineal $V-I$), mientras que los no óhmicos tienen relación no lineal (diodos, transistores).
    
    \item En conductores metálicos, la resistencia aumenta con la temperatura, mientras que en semiconductores generalmente disminuye.
\end{enumerate}

\section*{REFERENCIAS BIBLIOGRÁFICAS}
\begin{enumerate}
    \item Halliday, D., Resnick, R., \& Walker, J. (2014). \textit{Fundamentals of Physics}. Wiley.
    \item Nilsson, J. W., \& Riedel, S. A. (2015). \textit{Electric Circuits}. Pearson.
    \item Boylestad, R. L. (2012). \textit{Introductory Circuit Analysis}. Pearson.
    \item Serway, R. A., \& Jewett, J. W. (2018). \textit{Physics for Scientists and Engineers}. Cengage Learning.
\end{enumerate}

\vspace{1cm}

\begin{center}
    \rule{0.8\textwidth}{0.5pt}
    
    \textbf{Firma del Estudiante}
    
    \vspace{0.5cm}
    
    \textbf{Nombre:} \underline{\hspace{8cm}}
    
    \vspace{0.5cm}
    
    \textbf{Fecha:} \underline{\hspace{8cm}}
\end{center}

\end{document}